\documentclass[11pt,a4paper]{article}
\usepackage[T1]{fontenc}
\usepackage{microtype}
\usepackage{amsmath,amssymb,amsthm,mathtools}
\usepackage{mathrsfs}
\usepackage{aliascnt}
\usepackage{booktabs,array,longtable}
\usepackage{geometry}
\usepackage{xcolor}
\usepackage{hyperref}
\usepackage[nameinlink,noabbrev]{cleveref}

\geometry{top=2.45cm,bottom=2.45cm,left=2.75cm,right=2.75cm}
\setlength{\emergencystretch}{2em}
\raggedbottom
\hypersetup{
  colorlinks=true,
  linkcolor=blue!55!black,
  citecolor=blue!55!black,
  urlcolor=blue!55!black,
  pdfauthor={Akari Hayami}
}

\theoremstyle{plain}
\newtheorem{theorem}{Theorem}[section]
\newaliascnt{proposition}{theorem}
\newtheorem{proposition}[proposition]{Proposition}
\aliascntresetthe{proposition}
\newaliascnt{lemma}{theorem}
\newtheorem{lemma}[lemma]{Lemma}
\aliascntresetthe{lemma}
\newaliascnt{corollary}{theorem}
\newtheorem{corollary}[corollary]{Corollary}
\aliascntresetthe{corollary}
\theoremstyle{definition}
\newaliascnt{definition}{theorem}
\newtheorem{definition}[definition]{Definition}
\aliascntresetthe{definition}
\newaliascnt{example}{theorem}
\newtheorem{example}[example]{Example}
\aliascntresetthe{example}
\theoremstyle{remark}
\newaliascnt{remark}{theorem}
\newtheorem{remark}[remark]{Remark}
\aliascntresetthe{remark}
\crefname{proposition}{proposition}{propositions}
\Crefname{proposition}{Proposition}{Propositions}
\crefname{lemma}{lemma}{lemmas}
\Crefname{lemma}{Lemma}{Lemmas}
\crefname{corollary}{corollary}{corollaries}
\Crefname{corollary}{Corollary}{Corollaries}
\crefname{definition}{definition}{definitions}
\Crefname{definition}{Definition}{Definitions}
\crefname{example}{example}{examples}
\Crefname{example}{Example}{Examples}
\crefname{remark}{remark}{remarks}
\Crefname{remark}{Remark}{Remarks}
\numberwithin{equation}{section}

\newcommand{\N}{\mathbb N}
\newcommand{\R}{\mathbb R}
\newcommand{\Z}{\mathbb Z}
\newcommand{\C}{\mathbb C}
\newcommand{\T}{\mathbb T}
\newcommand{\F}{\mathbb F}
\newcommand{\kk}{\Bbbk}
\newcommand{\norm}[1]{\left\lVert#1\right\rVert}
\newcommand{\Aex}{\mathcal A_{\mathrm{ex}}}
\newcommand{\Hex}{H_{\mathrm{ex}}}
\newcommand{\Nex}{\mathcal N_{\mathrm{ex}}}
\DeclareMathOperator{\im}{im}
\DeclareMathOperator{\id}{id}
\DeclareMathOperator{\dist}{dist}
\DeclareMathOperator{\reach}{reach}
\DeclareMathOperator{\Fill}{Fill}
\DeclareMathOperator{\Fillw}{\widehat{Fill}}
\DeclareMathOperator{\End}{End}
\DeclareMathOperator{\supp}{supp}
\DeclareMathOperator{\rank}{rank}
\DeclareMathOperator{\ob}{ob}
\usepackage{graphicx}
\usepackage{tikz}
\usepackage{pgfplots}
\usepackage{caption}
\usepackage[section]{placeins}
\usetikzlibrary{arrows.meta,positioning,calc,fit,backgrounds}
\pgfplotsset{compat=1.18}
\definecolor{viewblue}{HTML}{245A81}
\definecolor{viewteal}{HTML}{16766D}
\definecolor{vieworange}{HTML}{AE591E}
\definecolor{viewgray}{HTML}{53616C}
\tikzset{
  >=Latex,
  every picture/.style={font=\small},
  vbox/.style={draw=viewblue,rounded corners=2pt,fill=viewblue!5,
    align=center,inner sep=7pt,line width=.7pt},
  vinput/.style={vbox,draw=vieworange,fill=vieworange!6},
  vresult/.style={vbox,draw=viewteal,fill=viewteal!6},
  varrow/.style={->,line width=.8pt,draw=viewblue},
  vconditional/.style={->,dashed,line width=.8pt,draw=vieworange},
  vlabel/.style={font=\footnotesize,align=center,fill=white,inner sep=2pt},
  vpoint/.style={circle,fill=viewblue,inner sep=2.5pt},
  vseed/.style={circle,draw=vieworange,fill=vieworange!20,inner sep=3pt}
}
\pgfplotsset{viewaxis/.style={
  axis lines=left,axis line style={viewgray},
  tick label style={font=\scriptsize},label style={font=\small},
  legend style={font=\scriptsize,draw=none,fill=white},
  grid=major,grid style={viewgray!12},
  every axis plot/.append style={line width=.85pt},
  scaled ticks=false
}}
\captionsetup[figure]{font=small,labelfont=bf,labelsep=period}
\renewcommand{\topfraction}{.9}
\renewcommand{\bottomfraction}{.85}
\renewcommand{\textfraction}{.08}
\renewcommand{\floatpagefraction}{.75}
\setcounter{topnumber}{3}
\setcounter{bottomnumber}{2}
\setcounter{totalnumber}{4}
\newcommand{\ViewHeading}[1]{\par\smallskip{\small\bfseries #1}\par\smallskip}
\hypersetup{pdftitle={Minimal Marked Response Modules for Relative Obstruction Data}}

\title{\textbf{Minimal Marked Response Modules}\\
\large For Relative Obstruction Data}
\author{Akari Hayami}
\date{Reorganized algebraic companion --- 16 September 2026}

\begin{document}
\maketitle

\begin{abstract}
Given a finite-dimensional quotient, an obstruction seed, and an ordered
family of endomorphisms preserving the boundary subspace, we construct
the smallest invariant subspace containing the seed image.  The
construction is the elementary generated-submodule closure, with a unital
algebra convention and a finite dimension bound stated explicitly.
For a fixed number of marks it is a functor to modules over the free
associative algebra.  Strict cochain transport is distinguished from
independence of the seed representative modulo boundaries.

Two finite-field examples separate a trivial descended action from
genuine one-step growth and show why the marks cannot be omitted.  A
calculation for the rank-one torus family explains a limitation of using
this construction as an arithmetic companion: its nontrivial relative
seed already spans a one-dimensional quotient.  Arithmetic maps between
different scales do not automatically provide endomorphisms of a fixed
quotient.  No character space, norm, recurrence spectrum, or geometric
offset is constructed from an unmarked complex.
\end{abstract}

\section{The supplied datum}
\label{sec:datum}

Fix a field $\kk$ and a nonnegative integer $t$.  All vector spaces are
finite dimensional.  The companion \emph{Foundations of Information
Exclusion and Relative Obstruction} explains how a relative connecting
class may supply a quotient seed.  The present note begins after that
choice, and after the additional endomorphism marks have been specified.
The elementary closure below is not presented as a new general theorem
in representation theory; its role is to fix the exact finite response
associated with the supplied data.

\begin{definition}[Finite marked obstruction datum]\label{def:datum}
A datum is
\[
 \mathfrak E=(\mathsf Z,\mathsf B,S,\widetilde\delta;G_1,\ldots,G_t),
\]
where $\mathsf B\subseteq\mathsf Z$ is a subspace, $S$ is a seed space,
$\widetilde\delta:S\to\mathsf Z$ is linear, and each
$G_j\in\End_\kk(\mathsf Z)$ satisfies $G_j(\mathsf B)\subseteq\mathsf B$.
The marks are an \emph{indexed tuple}: their order and labels are retained
and repetitions are permitted.  Define
\[
 \Hex(\mathfrak E)=\mathsf Z/\mathsf B,
 \qquad \ob_{\mathfrak E}(s)=[\widetilde\delta(s)],
 \qquad V_{\mathfrak E}=\im\ob_{\mathfrak E}.
\]
\end{definition}

\subsection{A relative-cochain realization}

\Cref{fig:response-datum} shows the seed map and the boundary-preserving
square defining each descended mark.
\begin{figure}[tbp]
\centering
\ViewHeading{(a) The supplied datum and its quotient seed}
\begin{tikzpicture}
\node[vinput] (seed) at (-4.8,0) {$S$};
\node[vbox] (cycles) at (0,0) {$\mathsf Z$};
\node[vresult] (quotient) at (4.8,0) {$H=\mathsf Z/\mathsf B$};
\draw[varrow] (seed)--(cycles) node[midway,above] {$\widetilde\delta$};
\draw[varrow] (cycles)--(quotient) node[midway,above] {$\pi$};
\draw[varrow] (seed.south) to[bend right=22] node[midway,below] {$\ob=\pi\widetilde\delta$} (quotient.south);
\node[font=\footnotesize] at (4.8,-1.55) {$V=\operatorname{im}\ob\subseteq H$};
\end{tikzpicture}
\ViewHeading{(b) Boundary preservation is what makes descent well defined}
\begin{tikzpicture}
\node[vbox] (upperleft) at (-3.3,0) {$\mathsf Z$};
\node[vbox] (upperright) at (3.3,0) {$\mathsf Z$};
\node[vresult] (lowerleft) at (-3.3,-1.7) {$H$};
\node[vresult] (lowerright) at (3.3,-1.7) {$H$};
\draw[varrow] (upperleft)--(upperright) node[midway,above] {$G_j$};
\draw[varrow] (lowerleft)--(lowerright) node[midway,below] {$\overline G_j$};
\draw[varrow] (upperleft)--(lowerleft) node[midway,left] {$\pi$};
\draw[varrow] (upperright)--(lowerright) node[midway,right] {$\pi$};
\node[vinput,text width=11cm] at (0,-3.2)
 {$z'=z+b$, $b\in\mathsf B$\quad$\Longrightarrow$\quad
  $G_jz'-G_jz=G_jb\in\mathsf B$\\
  hence $[G_jz']=[G_jz]$};
\end{tikzpicture}
\caption{The response starts with a specified seed map and an indexed
tuple of marks. In this figure $H=\Hex$. Relative exactness can supply
the quotient seed, but does not select the endomorphisms. The square
commutes precisely because the displayed marks preserve the boundary
subspace.}
\label{fig:response-datum}
\end{figure}

For a finite regular cellular pair $(K,A)$ with a finite-dimensional
cellular sheaf $\mathcal F$, or a finite CW pair with a finite-rank local
system $\mathcal F$, one may take
\[
 \mathsf Z=Z^1(K,A;\mathcal F),\qquad
 \mathsf B=B^1(K,A;\mathcal F),\qquad S\subseteq H^0(A;\mathcal F).
\]
Choose a linear degree-zero lift $E$ of $S$ to $K$ and put
$\widetilde\delta=d^0E$.  Then $\Hex=H^1(K,A;\mathcal F)$ and
$\ob$ is the restriction of the connecting map.  Exactness gives
\[
 \ob(s)\ne0\iff s\text{ has no global extension}.
\]
The cochain short exact sequence and this realization are explained in
the foundations companion; for local systems the background is the
relative theory with local coefficients \cite[Section 3.H]{hatcher}.
Boundary-preserving $G_j$ are not selected by that exact sequence.

\begin{proposition}[Independence of the quotient seed]
\label{prop:seed-choice}
Fix $\mathsf Z,\mathsf B,S$ and the marks.  If
$\widetilde\delta'-\widetilde\delta:S\to\mathsf Z$ takes values in
$\mathsf B$, then the two quotient obstruction maps are equal.  Every
construction below therefore gives the same response subspace.
\end{proposition}

\begin{proof}
The quotient map $\mathsf Z\to\mathsf Z/\mathsf B$ annihilates the
difference, so both seed images coincide.  The descended marks also
coincide because they were held fixed.
\end{proof}

In the relative realization, replacing the lift $E$ by $E'$ changes
$d^0E$ by a relative boundary, so satisfies the proposition.  This
independence statement does not require the identity maps to form a
strict morphism of the two displayed cochain data; strict morphisms below
require equality before taking the quotient.

\section{The unital algebra and finite closure}
\label{sec:closure}

Each mark descends to $\overline G_j\in\End_\kk(\Hex)$.  Let
\[
 \mathscr R_t=\kk\langle X_1,\ldots,X_t\rangle
\]
be the unital free associative algebra, with $\mathscr R_0=\kk$.
Its representation on $\Hex$ sends $X_j$ to $\overline G_j$ and
$1$ to $\id_{\Hex}$.  Define the effective algebra and response by
\begin{equation}\label{eq:response}
 \Aex(\mathfrak E)=\im(\mathscr R_t\to\End_\kk(\Hex)),
 \qquad \Nex(\mathfrak E)=\Aex(\mathfrak E)V_{\mathfrak E}.
\end{equation}
The product denotes the linear span of all such actions.  In particular,
the identity, equivalently the empty operator word, is included.  Thus
$V_{\mathfrak E}\subseteq\Nex(\mathfrak E)$ even with no marks.  On a
zero-dimensional quotient all actions, including the identity map, are
the unique zero map.

\begin{theorem}[Minimality, termination, and the dimension bound]
\label{thm:closure}
The response $\Nex(\mathfrak E)$ is the unique smallest subspace of
$\Hex$ containing $V_{\mathfrak E}$ and invariant under all descended
marks.  Define
\begin{equation}\label{eq:iteration}
 N^{(0)}=V_{\mathfrak E},\qquad
 N^{(r+1)}=N^{(r)}+\sum_{j=1}^t\overline G_jN^{(r)}.
\end{equation}
The iteration reaches $\Nex$ after at most
$\dim\Hex-\dim V_{\mathfrak E}$ strict dimension increases.  An equality
$N^{(r+1)}=N^{(r)}$ certifies termination at that step.
\end{theorem}

\begin{proof}
Induction identifies $N^{(r)}$ with the span of words of length at most
$r$ applied to $V_{\mathfrak E}$.  Their union is the generated module
in \cref{eq:response}.  Every invariant subspace containing the seed
contains each such word, giving minimality and uniqueness.

If equality occurs at one step, every $\overline G_jN^{(r)}$ already
lies in $N^{(r)}$, so all later steps are equal.  Before equality, each
step increases dimension by at least one.  Starting at
$\dim V_{\mathfrak E}$ and remaining within $\Hex$ proves the bound.
\end{proof}

If the seed is zero, the response is zero without any iteration.  If
$t=0$, the response is the seed image.  These cases are included in the
unital convention rather than treated by an implicit exception.

\subsection{A coordinate algorithm}

Choose a quotient basis and write $M_j$ for the matrices of the descended
marks.  If the columns of $Q_r$ form a basis for $N^{(r)}$, obtain $Q_{r+1}$
by taking a column-space basis of
\[
 [\,Q_r\mid M_1Q_r\mid\cdots\mid M_tQ_r\,].
\]
Stop when the rank no longer increases.  The initial columns come from
the seed map after quotienting by $\mathsf B$.  One need not enumerate
the whole effective algebra.  The dimension bound is a bound on strict
growth steps, not a bit-complexity estimate; field operations and exact
rank computation still have their own costs.

\Cref{fig:response-closure} depicts operator words and the exact
column-space iteration, with the empty word and stopping criterion
retained.
\begin{figure}[tbp]
\centering
\ViewHeading{(a) Words generate a linear subspace, not just an orbit set}
\begin{tikzpicture}
\node[vinput] (empty) at (0,0) {$V=1V$};
\node[vbox] (first) at (-3,-1.2) {$\overline G_1 V$};
\node[vbox] (second) at (3,-1.2) {$\overline G_2 V$};
\node (firstfirst) at (-4.8,-2.4) {$\overline G_1\overline G_1 V$};
\node (secondfirst) at (-1.6,-2.4) {$\overline G_2\overline G_1 V$};
\node (firstsecond) at (1.6,-2.4) {$\overline G_1\overline G_2 V$};
\node (secondsecond) at (4.8,-2.4) {$\overline G_2\overline G_2 V$};
\draw[varrow] (empty)--(first);\draw[varrow] (empty)--(second);
\draw[varrow] (first)--(firstfirst);\draw[varrow] (first)--(secondfirst);
\draw[varrow] (second)--(firstsecond);\draw[varrow] (second)--(secondsecond);
\node[font=\footnotesize] at (0,-3.15)
 {Take the sum of all displayed subspaces at every level; words can be linearly dependent.};
\end{tikzpicture}
\ViewHeading{(b) Exact column-space iteration}
\begin{tikzpicture}
\node[vinput,text width=4cm] (initial) at (-4.1,0) {$Q_0$: a basis of the quotient seed};
\node[vbox,text width=5.1cm] (step) at (2.8,0)
 {$Q_{m+1}=\operatorname{colbasis}$\\$[Q_m\mid M_1Q_m\mid\cdots\mid M_tQ_m]$};
\node[vinput,text width=5.1cm] (test) at (2.8,-1.65)
 {$\operatorname{rank}Q_{m+1}=\operatorname{rank}Q_m$?};
\node[vresult,text width=4cm] (end) at (-4.1,-1.65)
 {Stop: columns span $\Nex$};
\draw[varrow] (initial)--(step);
\draw[varrow] (step)--(test);
\draw[varrow] (test)--(end) node[midway,above,font=\footnotesize] {yes};
\draw[varrow] (test.east)--++(.8,0)|- (step.east)
 node[pos=.25,right,font=\footnotesize] {no: next $m$};
\node[font=\footnotesize,align=center] at (0,-2.9)
 {At most $\dim H-\dim V$ strict rank increases.\\
  $V=0\Rightarrow\Nex=0$; no marks $\Rightarrow\Nex=V$.};
\end{tikzpicture}
\caption{Panel (a) illustrates two indexed marks with the empty word
retained. It is a word diagram, not a claim of linear independence or
commutativity. Panel (b) works for arbitrary $t$ without enumerating the
effective algebra. The bound counts strict growth steps, not all field
operations or the final equality check.}
\label{fig:response-closure}
\end{figure}

\section{Strict maps and the functorial target}
\label{sec:functor}

\begin{definition}[Strict morphism]\label{def:strict}
Between data with the same indexed number $t$ of marks, a strict morphism
is a pair of linear maps
\[
 \Phi_{\mathsf Z}:\mathsf Z\to\mathsf Z',\qquad
 \Phi_S:S\to S'
\]
satisfying
\begin{align*}
 \Phi_{\mathsf Z}(\mathsf B)&\subseteq\mathsf B',\\
 \Phi_{\mathsf Z}\widetilde\delta&=\widetilde\delta'\Phi_S,\\
 \Phi_{\mathsf Z}G_j&=G'_j\Phi_{\mathsf Z}\qquad(1\leq j\leq t).
\end{align*}
\end{definition}

Identities and componentwise composition satisfy these conditions, so the
data form a category $\mathbf{Excl}_t$.  Every strict map induces a
quotient map $\overline\Phi:\Hex\to\Hex'$.

\begin{theorem}[Functoriality over a fixed free algebra]
\label{thm:functor}
The assignments $\mathfrak E\mapsto\Hex(\mathfrak E)$ and
$\mathfrak E\mapsto\Nex(\mathfrak E)$ define functors
\[
 \mathbf{Excl}_t\longrightarrow\mathscr R_t\text{-}\mathbf{Mod}_{\mathrm{fin}}.
\]
On response modules, the map is the restriction of $\overline\Phi$;
in particular,
\[
 \overline\Phi(\Nex(\mathfrak E))\subseteq\Nex(\mathfrak E').
\]
If $V_{\mathfrak E'}\subseteq\overline\Phi(\Nex(\mathfrak E))$, this
inclusion is an equality.  If in addition $\overline\Phi$ is injective
on $\Nex(\mathfrak E)$, its restriction is an isomorphism of responses.
\end{theorem}

\begin{proof}
Boundary preservation defines the quotient map.  The other conditions
descend to
\[
 \overline\Phi\,\ob_{\mathfrak E}
     =\ob_{\mathfrak E'}\Phi_S,\qquad
 \overline\Phi\,\overline G_j=\overline G'_j\,\overline\Phi.
\]
Thus $\overline\Phi$ is $\mathscr R_t$-linear and sends the seed image,
and every word applied to it, into the target response.  Identity and
composition laws are inherited from quotient maps.

The image of the source response is invariant under all target marks.
If it contains the target seed, target minimality gives the reverse
inclusion.  Injectivity on the response then gives the asserted
isomorphism.
\end{proof}

\Cref{fig:response-strict} separates the strict-morphism conditions from
the additional conditions for equality or isomorphism of responses.
\begin{figure}[tbp]
\centering
\begin{minipage}[t]{.48\linewidth}
\centering\ViewHeading{(a) Strict equality of seed representatives}
\begin{tikzpicture}
\node (seed) at (0,0) {$S$};
\node (cycles) at (3.7,0) {$\mathsf Z$};
\node (seednew) at (0,-1.7) {$S'$};
\node (cyclesnew) at (3.7,-1.7) {$\mathsf Z'$};
\draw[varrow] (seed)--(cycles) node[midway,above] {$\widetilde\delta$};
\draw[varrow] (seednew)--(cyclesnew) node[midway,below] {$\widetilde\delta'$};
\draw[varrow] (seed)--(seednew) node[midway,left] {$\Phi_S$};
\draw[varrow] (cycles)--(cyclesnew) node[midway,right] {$\Phi_{\mathsf Z}$};
\end{tikzpicture}
\end{minipage}\hfill
\begin{minipage}[t]{.48\linewidth}
\centering\ViewHeading{(b) Each indexed mark intertwines}
\begin{tikzpicture}
\node (left) at (0,0) {$\mathsf Z$};
\node (right) at (3.7,0) {$\mathsf Z$};
\node (leftnew) at (0,-1.7) {$\mathsf Z'$};
\node (rightnew) at (3.7,-1.7) {$\mathsf Z'$};
\draw[varrow] (left)--(right) node[midway,above] {$G_j$};
\draw[varrow] (leftnew)--(rightnew) node[midway,below] {$G'_j$};
\draw[varrow] (left)--(leftnew) node[midway,left] {$\Phi_{\mathsf Z}$};
\draw[varrow] (right)--(rightnew) node[midway,right] {$\Phi_{\mathsf Z}$};
\end{tikzpicture}
\end{minipage}
\medskip
\ViewHeading{(c) Boundary preservation and the induced response map}
\begin{tikzpicture}
\node[vinput,text width=11.8cm] (boundary) at (0,0)
 {$\Phi_{\mathsf Z}(\mathsf B)\subseteq\mathsf B'$ defines $\overline\Phi:H\to H'$};
\node[vresult,text width=11.8cm] (response) at (0,-1.45)
 {$\overline\Phi(\Nex)\subseteq\Nex'$ as $\mathscr R_t$-modules};
\draw[varrow] (boundary)--(response);
\node[font=\footnotesize,align=center,text width=12cm] at (0,-2.8)
 {Equality if the image contains the target seed $V'$.\\
  An isomorphism of responses if, in addition, $\overline\Phi|_{\Nex}$ is injective.};
\end{tikzpicture}
\caption{All three strict-morphism requirements are needed. Equality in
(a) holds before quotienting and is stronger than changing a seed
representative by a boundary. The number and labels of the marks remain
fixed. Identities and composition descend to the response restrictions.}
\label{fig:response-strict}
\end{figure}

\begin{remark}[Effective algebras need not have induced algebra maps]
\label{rem:effective}
The common functorial target is modules over $\mathscr R_t$, not modules
over one fixed effective algebra.  A strict map does not generally
induce an algebra homomorphism $\Aex(\mathfrak E)\to\Aex(\mathfrak E')$
sending every descended mark to its counterpart: operator relations may
hold on the source and fail on the target.

For example, over $\F_2$ with one mark, take a source quotient
$\F_2$ with mark zero and a target quotient $\F_2^2$ with mark
$\operatorname{diag}(0,1)$.  Give both data zero boundary subspaces,
seed space $\F_2$, and seed maps $s\mapsto s$ and $s\mapsto(s,0)$.
The inclusion $s\mapsto(s,0)$, together with the identity seed map, is
strict.  But a generator-respecting algebra map would send the zero
source mark to the nonzero target mark, which is impossible.  The
$\mathscr R_1$-module map is nevertheless well defined.
\end{remark}

\Cref{fig:response-algebras} displays the preceding counterexample to an
induced generator-respecting map of effective algebras.
\begin{figure}[tbp]
\centering
\begin{tikzpicture}
\node[vbox] (free) at (0,0) {$\mathscr R_1=\F_2\langle X\rangle$};
\node[vbox,text width=5.6cm] (source) at (-3.55,-1.8)
 {Source $H=\F_2$\\$X\mapsto0$\\$\Aex=\F_2\id$};
\node[vbox,text width=5.6cm] (target) at (3.55,-1.8)
 {Target $H'=\F_2^2$\\$X\mapsto\operatorname{diag}(0,1)$\\$\Aex'=\langle\id,\operatorname{diag}(0,1)\rangle$};
\draw[varrow] (free)--(source);
\draw[varrow] (free)--(target);
\node[vresult,text width=11.9cm] (module) at (0,-3.85)
 {$\overline\Phi(s)=(s,0)$ is $\mathscr R_1$-linear and carries the seed to the seed};
\node[vinput,text width=11.9cm] (failure) at (0,-5.25)
 {No generator-respecting map $\Aex\to\Aex'$:\\
  it would have to send $0$ to $\operatorname{diag}(0,1)\neq0$.};
\end{tikzpicture}
\caption{A common free algebra acts on both quotients, but their
effective images impose different relations. Both boundary spaces are
zero, the seed maps are $s\mapsto s$ and $s\mapsto(s,0)$, and the seed
comparison is the identity. The strict module map exists even though the
claimed generator-respecting algebra map does not.}
\label{fig:response-algebras}
\end{figure}

\section{Finite witnesses and dependence on marks}
\label{sec:examples}

\begin{example}[A descended identity]\label{ex:identity}
Over $\F_2$, set
\[
 \mathsf Z=\F_2^2,\quad\mathsf B=\langle(1,1)\rangle,\quad
 S=\F_2,\quad\widetilde\delta(s)=(s,0),\quad G(x,y)=(y,x).
\]
The mark preserves $\mathsf B$.  The quotient is one dimensional, and
$[(1,0)]=[(0,1)]\ne0$ because their difference is $(1,1)$.
Thus $\overline G=\id$, the seed already spans the quotient, and
\[
 \Aex=\F_2\id,\qquad N^{(0)}=\Nex=\Hex.
\]
This example verifies descent and a nonzero quotient seed, but it does
not exhibit response growth.
\end{example}

\begin{example}[Strict growth and a nontrivial boundary quotient]
\label{ex:growth}
Let $\mathsf Z=\F_2^3$ with its ordered coordinate basis,
$\mathsf B=\langle(0,0,1)\rangle$, $S=\F_2$, and
$\widetilde\delta(s)=(s,0,0)$.  Choose
\[
 G(x,y,z)=(0,x,0).
\]
The mark annihilates $\mathsf B$ and therefore descends.  In the quotient
basis $[(1,0,0)],[(0,1,0)]$, it has matrix
\[
 \overline G=\begin{pmatrix}0&0\\1&0\end{pmatrix},\qquad
 \overline G^2=0.
\]
The iteration gives
\[
 N^{(0)}=\langle[(1,0,0)]\rangle
 \subsetneq N^{(1)}=\Hex=N^{(2)}.
\]
Moreover, $\Aex=\langle\id,\overline G\rangle$ is a two-dimensional
algebra, isomorphic to $\F_2[X]/(X^2)$.  The single strict growth step
attains the bound $\dim\Hex-\dim N^{(0)}=1$.
\end{example}

\begin{proposition}[The unmarked seed does not determine the response]
\label{prop:marks-needed}
Keeping $\mathsf Z,\mathsf B,S,\widetilde\delta$ fixed does not in
general determine the dimension of the response.
\end{proposition}

\begin{proof}
Use the underlying data of \cref{ex:growth}.  Its displayed mark gives
a two-dimensional response.  Replacing that mark by the zero map leaves
the unmarked quotient and seed unchanged, but every iterate equals the
one-dimensional seed image.  The identity on the underlying spaces is
not a strict morphism intertwining these different marks.
\end{proof}

\Cref{fig:response-witnesses} gives finite-point views of the examples,
including the dependence of the response on the chosen mark.
\begin{figure}[tbp]
\centering
\begin{minipage}[t]{.47\linewidth}
\centering\ViewHeading{(a) Nilpotent mark: strict growth}
\begin{tikzpicture}
\node[vseed,label=below:$0$] (zero) at (0,0) {};
\node[vseed,label=below:{$[e_1]$}] (first) at (2.8,0) {};
\node[vpoint,label=above:{$[e_2]$}] (second) at (0,1.6) {};
\node[vpoint,label=above:{$[e_1+e_2]$}] (sum) at (2.8,1.6) {};
\draw[varrow] (first)--(second) node[midway,above right,font=\footnotesize] {$\overline G$};
\draw[varrow] (second)--(zero);
\draw[varrow] (sum)--(second);
\node[align=center] at (1.4,-1.35)
 {$N^{(0)}=\{0,[e_1]\}$\\$N^{(1)}=H=\F_2^2$};
\end{tikzpicture}
\end{minipage}\hfill
\begin{minipage}[t]{.47\linewidth}
\centering\ViewHeading{(b) Zero mark: no growth}
\begin{tikzpicture}
\node[vseed,label=below:$0$] (zero) at (0,0) {};
\node[vseed,label=below:{$[e_1]$}] (first) at (2.8,0) {};
\node[circle,draw=viewgray,inner sep=2.5pt,label=above:{$[e_2]$}] (second) at (0,1.6) {};
\node[circle,draw=viewgray,inner sep=2.5pt,label=above:{$[e_1+e_2]$}] (sum) at (2.8,1.6) {};
\draw[varrow] (first)--(zero);
\draw[varrow] (second)--(zero);
\draw[varrow] (sum)--(zero);
\node[align=center] at (1.4,-1.35)
 {$N^{(m)}=\{0,[e_1]\}$ for all $m$\\$\dim\Nex=1$};
\end{tikzpicture}
\end{minipage}
\ViewHeading{(c) A swap can descend to the identity}
\begin{tikzpicture}
\node[vbox,text width=4.8cm] (swap) at (-3.8,0)
 {$\mathsf Z=\F_2^2$, $G(x,y)=(y,x)$\\$(1,0)\longleftrightarrow(0,1)$};
\node[vresult,text width=5.3cm] (identity) at (3.8,0)
 {$\mathsf B=\langle(1,1)\rangle$\\$[(1,0)]=[(0,1)]$,\quad $\overline G=\id$};
\draw[varrow] (swap)--(identity) node[midway,above,font=\footnotesize] {quotient};
\end{tikzpicture}
\caption{Exact finite-field witnesses, drawn as finite points rather
than continuous planes. Panels (a) and (b) use the same quotient of
$\F_2^3$ by $\langle e_3\rangle$ and the same seed; only the mark changes.
Orange points are the initial seed subspace, blue points are added by
closure, and hollow points are outside the response. Zero is fixed in
both panels (its self-loop is omitted). In (a), linear span adds
$[e_1+e_2]$ as well as the new direction. Panel (c) already has $V=H$.}
\label{fig:response-witnesses}
\end{figure}

\section{What the arithmetic torus does and does not supply}
\label{sec:torus}

For the one-vertex torus $K_d$ and local system $\mathcal L_x$ in
\emph{Arithmetic Local-System Filling Spectra}, write
$v_j=e^{2\pi\mathrm i x_j}-1$ and $v=(v_1,\ldots,v_d)$.
Relative to the marked product cells, the cochain differentials begin
\[
 d_x^0(s)=vs,\qquad
 (d_x^1w)_{jk}=v_jw_k-v_kw_j\quad(j<k).
\]
For $d=1$ there are no degree-two coordinates.

\begin{proposition}[The nontrivial torus seed already fills its quotient]
\label{prop:torus}
If $x\ne0$ in $\T^d$, then
\[
 H^1(K_d,\{a\};\mathcal L_x)=\langle[v]\rangle
\]
is one dimensional and the connecting seed of the vertex state $1$
spans it.  Consequently every choice of admissible endomorphism marks
on this relative datum has $\Nex=\Hex$, with no strict growth steps.
At $x=0$ the relative quotient has dimension $d$, but the connecting
seed is zero and its response is zero for every choice of marks.
\end{proposition}

\begin{proof}
Relative degree-zero cochains vanish, so $B^1(K_d,\{a\};\mathcal L_x)=0$.
For $v\ne0$, choose a coordinate $j$ with $v_j\ne0$.  The cocycle
equations give $w_k=v_k w_j/v_j$ for every $k$, so
$\ker d_x^1=\langle v\rangle$.  This also holds for $d=1$, when the
one-dimensional degree-one space is spanned by the nonzero vector $v$.
The connecting cochain is $v$, proving the seed assertion and hence the
response assertion by \cref{thm:closure}.

At $x=0$, both displayed differentials vanish, so relative $H^1=\C^d$
while the connecting seed is zero.  The zero seed generates zero.
\end{proof}

The wedge used for the chain filling calculation and the full torus
used for relative cohomology share their degree-zero coboundary, but not
their higher-degree cochain groups.  In particular, the wedge has no
degree-two cocycle equations, so its relative degree-one quotient must
not be substituted into \cref{prop:torus} without changing the datum.

The arithmetic maps $T_{q,k}$ are maps between chain complexes at different
scales with different powered coefficients.  They are not, as given,
endomorphisms of the fixed relative quotient used above.  To turn them
into response marks one must specify appropriate quotient identifications
or a new finite model and prove boundary preservation and intertwining.
Trivial powered characters also change the torus relative dimension, so
uniform identifications cannot simply be assumed.

Thus a nontrivial response-growth example is not automatically provided
by the one-vertex torus calculation.  Larger quotients and meaningful
marks require additional input, such as the deliberately supplied data
in \cref{ex:growth}.

\Cref{fig:response-torus} summarizes the torus saturation calculation
and the missing comparison between arithmetic scale maps and response
endomorphisms.
\begin{figure}[tbp]
\centering
\ViewHeading{(a) The full-torus seed leaves no room for strict growth}
\begin{tabular}{@{}llll@{}}
\toprule
Parameter & Relative quotient $H$ & Seed $V$ & Every admissible set of marks\\\midrule
$x\neq0$ & $\langle[v]\rangle$, dimension $1$ & $H$ & $\Nex=H$, no growth\\[3pt]
$x=0$ & $\C^d$, dimension $d$ & $0$ & $\Nex=0$, no growth\\\bottomrule
\end{tabular}
\medskip
\ViewHeading{(b) Different types of maps}
\begin{tikzpicture}
\node[vbox,text width=5.6cm] (scales) at (-3.55,0)
 {\textbf{Arithmetic scale map}\\
  $T_{q,k}:\mathcal C_x(kq)_1\to\mathcal C_x(q)_1$\\
  Different scales and powered coefficients};
\node[vresult,text width=5.6cm] (marks) at (3.55,0)
 {\textbf{Response endomorphism}\\$\overline G_j:H\to H$\\One supplied fixed quotient};
\node[vinput,text width=11.9cm] (extra) at (0,-2)
 {Missing additional construction: quotient identifications or a new finite model,\\
  with proved boundary preservation and intertwining};
\draw[vconditional] (scales.south)--(extra.north west);
\draw[vconditional] (extra.north east)--(marks.south);
\node[font=\footnotesize,align=center] at (0,-3.35)
 {Dashed arrows are requirements, not maps constructed by the papers.\\
  Trivial powered characters can also change the relative dimension.};
\end{tikzpicture}
\caption{The arithmetic torus does not automatically provide a growing
response module. Its full degree-one cocycle equations force the top
table; substituting the wedge changes the datum. The lower comparison
keeps cross-scale chain maps separate from fixed-quotient marks.}
\label{fig:response-torus}
\end{figure}

\section{Scope and verification}

This note assigns a minimal response to a supplied marked finite datum.
It does not choose the marks, make them canonical under unmarked homotopy,
extract a finite model from an infinite family, or construct a parameter
metric, Diophantine spectrum, or offset geometry.  Its functoriality is
exactly the strict transport in \cref{def:strict}, while its invariance
under changes of seed representative is the separate quotient statement
in \cref{prop:seed-choice}.

The shared verifier \path{scripts/verify_restructured_series.py}
enumerates the finite-field examples, checks boundary preservation and
the dimensions of their iterated quotients, and checks the displayed
strict-map example.  The finite calculations illustrate the theorem;
termination and functoriality follow from the written proofs.

\bibliographystyle{plain}
\bibliography{series_references}
\end{document}